\section{电荷、电压、电流与参考网络}

物质由带电粒子组成。电荷是物理属性，电流描述电荷通过截面的速率：

\[
I=\frac{\mathrm{d}Q}{\mathrm{d}t}.
\]

\(\SI{1}{A}=\SI{1}{C/s}\)。传统电流方向定义为正电荷移动方向；金属导线中
主要载流电子的平均漂移方向与传统电流相反。分析电路时保持传统电流方向即可，
不需要每次改用电子运动方向。

\term{电压}{两点之间的电势差}表示单位电荷可交换的能量：

\[
V=\frac{W}{Q}.
\]

因此电压必须在两个点之间测量。“某引脚是 \(\SI{3.3}{V}\)”通常省略了
“相对于该电路的参考地 GND”。两个没有共同参考的系统即使都标着
\(\SI{3.3}{V}\)，也不能直接假定信号兼容。

\begin{fromzerobox}
导线是低阻连接，不是理想的“瞬间传送”。真实导线有电阻、电感和对周围导体
的电容；频率越高、边沿越快，几何形状和回流路径越重要。低速原理图可以先把
短导线视为同一节点，高速章节再恢复被忽略的寄生参数。
\end{fromzerobox}

\topic{从原子到导线：电荷为什么能移动}
原子由带正电的原子核和带负电的电子组成。宏观物体通常近似电中性，但电子的
分布和移动会产生电场、电势差与电流。金属中有较容易响应电场的载流电子，
绝缘体中的电荷则更难长距离移动；半导体的导电能力可以通过材料、掺杂、电场、
光和温度控制，因此适合制造二极管和晶体管。

这里不需要先掌握量子力学，但要保留三个边界：

\begin{itemize}[leftmargin=2.2em]
  \item 电流不是导线里储存的一桶“电”；它是电荷流动的速率；
  \item 信号变化由电磁场沿互连传播，不等到某个电子从电源跑完整条导线；
  \item 导体、绝缘体和半导体不是绝对分类，温度、电压和结构会改变行为。
\end{itemize}

\topic{库仑、电流和“每秒有多少电子”}
一个电子携带的电荷量大小约为

\[
e=1.602\times10^{-19}\ \mathrm{C}.
\]

若电流为 \(I\)，每秒通过截面的电子数量量级为

\[
N=\frac{I}{e}.
\]

\begin{workedexample}
\(\SI{1}{mA}\) 电流每秒对应的电子数量约为

\[
N=\frac{1\times10^{-3}}{1.602\times10^{-19}}
\approx6.24\times10^{15}.
\]

这个巨大数量解释了为什么宏观电流看起来连续，但数字电路最终仍由微观载流子
行为实现。计算不表示这些电子都从电源出发后在一秒内穿过整块主板。
\end{workedexample}

\topic{电压是两点之差，GND 是人为选择的参考网络}
电势本身取决于参考；电压是两点电势之差：

\[
V_{AB}=V_A-V_B,\qquad V_{BA}=-V_{AB}.
\]

把某网络命名为 GND，通常表示设计者选择它作为本电路的零电势参考和主要回流
网络。它不自动等于大地，也不保证板上任意两个“地”点在高速或大电流条件下
严格没有电压差。

\begin{osgridlongtable}{L{0.18\linewidth}L{0.33\linewidth}L{0.39\linewidth}}
名称 & 常见职责 & 不能自动推出 \\ \hline
\endfirsthead
名称 & 常见职责 & 不能自动推出 \\ \hline
\endhead
电路 GND & 信号与电源的公共参考/回流网络 & 必然与大地相连 \\ \hline
保护地 PE & 故障电流安全路径和外壳保护 & 可以随意当作高速信号回流 \\ \hline
机壳地/chassis & 屏蔽、连接器外壳和 ESD 泄放 & 与数字地处处短接 \\ \hline
模拟地/数字地 & 表达噪声与回流分区意图 & 原理图中名字不同就完全隔离 \\ \hline
隔离侧地 & 隔离器某一侧的局部参考 & 与另一侧可直接连普通信号 \\ \hline
\end{osgridlongtable}

\begin{center}
\begin{tikzpicture}[font=\footnotesize\sffamily,>=Latex]
  \node[draw=BookBlue,fill=BookCodeBg,minimum width=30mm,minimum height=14mm,
        align=center]
    (a) at (0,1.8) {系统 A\\\(V_A=\SI{3.3}{V}\) 相对 \(GND_A\)};
  \node[draw=BookTeal,fill=BookCodeBg,minimum width=30mm,minimum height=14mm,
        align=center]
    (b) at (7.2,1.8) {系统 B\\\(V_B=\SI{3.3}{V}\) 相对 \(GND_B\)};
  \node[os state] (ga) at (0,0) {\(GND_A\)};
  \node[os state] (gb) at (7.2,0) {\(GND_B\)};
  \draw[os arrow] (ga) -- (a);
  \draw[os arrow] (gb) -- (b);
  \draw[BookRed,very thick,dashed] (a) -- node[above,align=center]
    {只连信号，\\共同参考未知} (b);
  \node[BookRed,align=center] at (3.6,-0.8)
    {\(V_A\) 与 \(V_B\) 数字相同，\\不代表接收端看到的差分电压相同};
\end{tikzpicture}
\end{center}

\topic{电源不是“恒定电压标签”，而是带能力边界的能量源}
理想电压源无论负载电流多少都保持电压；真实电源具有最大电流、内部阻抗、
动态响应、纹波、启动时间、保护与温升。一个简单等效模型是理想电压
\(V_s\) 串联内部电阻 \(R_s\)：

\[
V_{\mathrm{load}}=V_s-I_{\mathrm{load}}R_s.
\]

\begin{center}
\begin{circuitikz}[american voltages]
  \draw (0,0) to[battery1,l_={理想源 \(V_s\)}] (0,3)
    to[R,l={内部阻抗 \(R_s\)}] (3,3)
    -- (4,3)
    to[R,l={负载 \(R_L\)}] (4,0)
    -- (0,0);
  \draw[->,BookBlue,thick] (1.1,3.45) -- (2.4,3.45)
    node[midway,above]{\(I_{\mathrm{load}}\)};
  \draw[<->,BookTeal,thick] (4.7,0.1) -- (4.7,2.9)
    node[midway,right]{\(V_{\mathrm{load}}\)};
\end{circuitikz}
\end{center}

\begin{workedexample}
\(\SI{5}{V}\) 电源等效内部电阻 \(\SI{0.2}{\ohm}\)。负载从
\(\SI{0.1}{A}\) 突增到 \(\SI{1.5}{A}\)，仅按静态模型估算：

\[
V_{\mathrm{load}}=\SI{5}{V}
-\SI{1.5}{A}\times\SI{0.2}{\ohm}
=\SI{4.7}{V}.
\]

若芯片最低允许电压是 \(\SI{4.75}{V}\)，静态就已越界。真实瞬态还要叠加
电源控制环路响应、线缆/PCB 电感和去耦电容，因此需要示波器而不是只看电源
铭牌。
\end{workedexample}

\topic{恒压、恒流与限流状态}
台式电源设定 \(\SI{5}{V}\)、限流 \(\SI{0.5}{A}\) 时，负载需求低于
\(\SI{0.5}{A}\) 通常工作在恒压模式；若负载试图取得更大电流，电源会降低
输出电压以把电流限制在设定值。此时看到板上电压很低，不一定是电源损坏，
可能是短路或负载过重触发限流。

\topic{节点方程应先于“电流从哪条画线走”}
对任一节点选择流入为正、流出为负：

\[
\sum_k I_k=0.
\]

方向可以先任意假设；计算结果为负表示真实方向与假设相反。关键是整组方程
保持同一符号约定。

\begin{center}
\begin{circuitikz}[american currents]
  \draw (0,3) node[vcc]{\(\SI{5}{V}\)}
    to[R,l={\(R_1=\SI{1}{k\ohm}\)},i>^={\(I_1\)}] (0,1.5)
    -- (2.2,1.5) node[circ,label=above:{节点 \(X\)}] (x) {};
  \draw (x) to[R,l={\(R_2=\SI{2}{k\ohm}\)},i>^={\(I_2\)}] (2.2,0)
    node[ground]{};
  \draw (x) -- (4.4,1.5)
    to[R,l={\(R_3=\SI{3}{k\ohm}\)},i>^={\(I_3\)}] (4.4,0)
    node[ground]{};
\end{circuitikz}
\end{center}

若节点电压为 \(V_X\)，则

\[
\frac{\SI{5}{V}-V_X}{R_1}
=\frac{V_X}{R_2}+\frac{V_X}{R_3}.
\]

这个方法可扩展到复杂偏置、电源监控和模拟接口。图画得弯曲还是笔直不影响
方程，节点连接才影响。

\topic{戴维南等效把复杂网络缩成一个源和一个电阻}
从两个端口向内看，许多线性网络可等效为开路电压
\(V_{\mathrm{th}}\) 串联等效电阻 \(R_{\mathrm{th}}\)。它不会保留内部每条支路，
却能回答接入不同负载时端口电压怎样变化。

以分压器为例：

\[
V_{\mathrm{th}}=V_{\mathrm{in}}\frac{R_2}{R_1+R_2},
\qquad
R_{\mathrm{th}}=R_1\parallel R_2.
\]

这比记住“分压器不能带负载”更精确：负载与
\(R_{\mathrm{th}}\) 相比足够大时，误差才足够小。

\topic{直流、交流、周期信号和数字边沿}
\term{直流}{DC}在所关注时间尺度上方向和大小近似稳定；
\term{交流}{AC}随时间变化并可能改变方向。数字信号虽然只被解释为有限状态，
其引脚电压仍是随时间变化的模拟波形，包含上升时间、下降时间、过冲、振铃和
噪声。

正弦信号可写为

\[
v(t)=V_{\mathrm{pk}}\sin(2\pi f t+\phi).
\]

其中 \(V_{\mathrm{pk}}\) 是峰值，\(f\) 是频率，\(\phi\) 是相位。对正弦波：

\[
V_{\mathrm{pp}}=2V_{\mathrm{pk}},
\qquad
V_{\mathrm{rms}}=\frac{V_{\mathrm{pk}}}{\sqrt{2}}.
\]

峰值、峰峰值和 RMS 回答不同问题；仪器读数若未说明是哪一种，就可能差一倍或
\(\sqrt{2}\)。

\topic{电容和电感对变化速度敏感}
电容与电感的理想关系：

\[
i_C=C\frac{\mathrm{d}v_C}{\mathrm{d}t},
\qquad
v_L=L\frac{\mathrm{d}i_L}{\mathrm{d}t}.
\]

电容反对电压瞬变，电感反对电流瞬变。对频率为 \(f\) 的正弦稳态，阻抗大小
分别为

\[
|X_C|=\frac{1}{2\pi f C},
\qquad
|X_L|=2\pi f L.
\]

因此同一颗去耦电容在低频和高频下表现不同；真实寄生电感还会让它在足够高频
不再像理想电容。

\begin{workedexample}
\(C=\SI{100}{nF}\)。在 \(\SI{1}{kHz}\) 时：

\[
|X_C|\approx\frac{1}{2\pi\times10^3\times100\times10^{-9}}
\approx\SI{1.59}{k\ohm}.
\]

在 \(\SI{10}{MHz}\) 时理想值约为 \(\SI{0.159}{\ohm}\)。但实际电容的封装和
走线电感会限制高频效果，所以“容量越大越好”并不成立，布局与频率范围同样
重要。
\end{workedexample}

\begin{failurebox}
用直流万用表读到 \(\SI{3.3}{V}\)，只能说明它在仪表响应范围内的平均/稳定
分量大致正常，不能证明没有短时跌落、过冲或噪声。判断复位、电源完整性和时钟
必须选择足够带宽、正确接地的测量方式。
\end{failurebox}


\section{电路定律、等效模型与动态响应}

电流要形成闭合回路。电源把能量送入回路，负载把电能转成光、热、机械运动
或内部状态变化。下图中电流从电源正端经过限流电阻和发光二极管，再沿地返回
电源负端。

\begin{center}
\begin{circuitikz}[american voltages]
  \draw
    (0,0) to[battery1,l_={\(\SI{5}{V}\)}] (0,3)
    to[short] (1.2,3)
    to[R,l={\(R_1\ \SI{330}{\ohm}\)}] (3.6,3)
    to[leD,l={\(D_1\)}] (5.3,3)
    to[short] (5.3,0)
    -- (0,0);
  \draw[->,BookBlue,thick] (1.2,3.35) -- (2.5,3.35) node[midway,above]{传统电流 \(I\)};
  \node[ground] at (2.65,0) {};
\end{circuitikz}
\end{center}
\diagramnote{元件符号表达电气关系，不表达实物尺寸或 PCB 位置。\(R_1\) 和
\(D_1\) 是位号；\(\SI{330}{\ohm}\) 是属性；导线交点和网络标签决定连接。}

把所有由理想导线相连的点看作一个\term{节点}{同一电位网络}。元件连接在
节点之间形成支路。读原理图时先找节点，不要按纸面上“线从左向右流动”的外观
猜测电气关系。

\topic{欧姆定律与功率}
对线性电阻，电压、电流和电阻满足：

\[
V=IR,\qquad I=\frac{V}{R},\qquad R=\frac{V}{I}.
\]

功率为

\[
P=VI=I^2R=\frac{V^2}{R}.
\]

这些等式中的方向与正负号必须一致；入门计算可以先使用大小，再检查元件额定
功率和温升。

\begin{workedexample}
用 \(\SI{3.3}{V}\) GPIO 驱动红色 LED，假定 LED 在目标电流附近压降约
\(\SI{1.9}{V}\)，期望电流 \(\SI{4}{mA}\)。串联电阻近似为

\[
R=\frac{\SI{3.3}{V}-\SI{1.9}{V}}{\SI{4}{mA}}
=\SI{350}{\ohm}.
\]

可选择邻近标准值 \(\SI{360}{\ohm}\)。此时

\[
I\approx\frac{\SI{1.4}{V}}{\SI{360}{\ohm}}
=\SI{3.89}{mA},
\qquad
P_R=I^2R\approx\SI{5.4}{mW}.
\]

还必须检查 GPIO 最大输出电流、输出高电压在该电流下是否仍满足规格、LED
极性与脉冲条件。公式给出候选值，数据手册决定候选值是否安全。
\end{workedexample}

\topic{基尔霍夫定律：电流不能凭空消失，能量不能凭空出现}
对一个节点，流入电流代数和等于流出电流代数和：

\[
\sum I_k=0.
\]

对一个闭合回路，各元件电压升降代数和为零：

\[
\sum V_k=0.
\]

它们分别来自电荷守恒和能量守恒。复杂电路仍可拆成节点方程和回路方程。

\topic{串联和并联}
串联电阻流过相同电流：

\[
R_{\mathrm{series}}=R_1+R_2+\cdots.
\]

并联电阻两端电压相同：

\[
\frac{1}{R_{\mathrm{parallel}}}
=\frac{1}{R_1}+\frac{1}{R_2}+\cdots.
\]

两个并联电阻可写成

\[
R_{\mathrm{parallel}}=\frac{R_1R_2}{R_1+R_2}.
\]

\section{无源元件、半导体器件与基础电路}

电阻不是“把多余电压吃掉”的黑盒。它在给定两端电压下限制电流，并把电能
转化为热。除阻值外，还要关注误差、额定功率、温度系数和封装。

\topic{分压器}
\begin{center}
\begin{circuitikz}[american voltages]
  \draw (0,3) node[vcc]{\(V_{\mathrm{in}}\)}
    to[R,l={\(R_1\)}] (0,1.5)
    to[R,l={\(R_2\)}] (0,0) node[ground]{};
  \draw (0,1.5) -- (2,1.5) node[right]{\(V_{\mathrm{out}}\)};
\end{circuitikz}
\end{center}

若输出端没有额外负载，则

\[
V_{\mathrm{out}}
=V_{\mathrm{in}}\frac{R_2}{R_1+R_2}.
\]

\begin{workedexample}
\(V_{\mathrm{in}}=\SI{5}{V}\)，\(R_1=\SI{10}{k\ohm}\)，
\(R_2=\SI{20}{k\ohm}\)：

\[
V_{\mathrm{out}}=\SI{5}{V}\times\frac{20}{10+20}
\approx\SI{3.33}{V}.
\]

但这不能直接当作大电流 \(\SI{3.3}{V}\) 电源。输出接入负载 \(R_L\) 后，
下臂变成 \(R_2\parallel R_L\)，输出电压会下降。分压器适合偏置和测量，
稳定供电通常需要稳压器。
\end{workedexample}

\topic{上拉、下拉与悬空输入}
数字输入阻抗通常很高。若没有器件主动驱动，引脚会受泄漏、耦合和静电影响，
读到不确定状态。上拉电阻让释放状态偏向电源，下拉电阻让释放状态偏向地。

\begin{center}
\begin{circuitikz}
  \draw (0,3) node[vcc]{\(\SI{3.3}{V}\)}
    to[R,l={\(R_{\mathrm{pull}}\ \SI{10}{k\ohm}\)}] (0,1.5)
    -- (2,1.5) node[right]{GPIO 输入};
  \draw (0,1.5) to[push button,l_={按键}] (0,0) node[ground]{};
\end{circuitikz}
\end{center}

按键松开时电阻把输入拉高；按下时输入接地。按下电流约为
\(\SI{3.3}{V}/\SI{10}{k\ohm}=\SI{0.33}{mA}\)。软件看到的是松开 1、
按下 0，称为\term{低有效}{asserted 时为低电平}。

\topic{电容：储存电场能量并反对电压突变}
电容的基本关系是

\[
Q=CV,\qquad i=C\frac{\mathrm{d}v}{\mathrm{d}t},
\qquad E=\frac{1}{2}CV^2.
\]

电容两端电压不能在有限电流下瞬间跳变。它可用于去耦、滤波、定时、交流耦合
和能量缓冲。真实电容还有等效串联电阻、等效串联电感、漏电、容差和耐压。

\topic{RC 充电与时间常数}
电阻 \(R\) 给电容 \(C\) 充电时：

\[
V_C(t)=V_{\mathrm{DD}}\left(1-e^{-t/(RC)}\right),
\qquad \tau=RC.
\]

在 \(t=\tau\) 时约达到最终值的 \(63.2\%\)，\(5\tau\) 时超过 \(99\%\)。

\begin{center}
\begin{circuitikz}
  \draw (0,2.5) node[vcc]{\(V_{\mathrm{DD}}\)}
    to[R,l={\(R\)}] (2,2.5)
    -- (3.5,2.5) node[right]{RESET\_N};
  \draw (2,2.5) to[C,l_={\(C\)}] (2,0) node[ground]{};
\end{circuitikz}
\end{center}

简单 RC 可以产生延迟，但真实复位通常还要保证明确阈值、迟滞、掉电快速复位、
电源斜率范围和多个电源域顺序，因此常使用专用 supervisor 或 PMIC 输出，
不能看到 RC 就假定复位一定可靠。

\begin{workedexample}
\(R=\SI{100}{k\ohm}\)，\(C=\SI{100}{nF}\)：

\[
\tau=RC=100\,000\times100\times10^{-9}
=\SI{10}{ms}.
\]

若输入高阈值假设为电源的 \(70\%\)，求过阈时间：

\[
0.7=1-e^{-t/\tau}
\Rightarrow
t=-\tau\ln(0.3)
\approx1.204\tau
\approx\SI{12.0}{ms}.
\]

这个结果依赖阈值、容差和电源斜率；不能把 \(\tau\) 直接等同于复位保持时间。
\end{workedexample}

\topic{去耦电容为什么要靠近电源引脚}
芯片内部大量晶体管同时翻转会在很短时间内需要电流。远处电源经导线和 PCB
平面的寄生电感来不及提供快速变化电流，局部电压可能下降。靠近电源脚的去耦
电容提供短时电流，并让高频回流形成小回路。

\begin{pdfdiagram}
  \node[os hardware] (reg) {稳压器\\较慢能量来源};
  \node[os hardware,right=25mm of reg] (cap) {去耦电容\\局部储能};
  \node[os hardware,right=25mm of cap] (ic) {芯片 VDD/GND\\瞬态负载};
  \draw[os arrow] (reg) -- node[above]{补充平均能量} (cap);
  \draw[os arrow] (cap) -- node[above]{短小高频回路} (ic);
  \draw[os arrow,bend left=35] (ic.south) to node[below]{回流} (cap.south);
\end{pdfdiagram}

\topic{二极管与 LED：方向和压降不能忽略}
二极管主要允许一个方向导通。硅 PN 结常被入门材料近似为正向压降约
\(\SI{0.7}{V}\)，但实际压降随器件、温度和电流变化；肖特基、LED、稳压管
又有不同用途和曲线。

\begin{center}
\begin{circuitikz}
  \draw (0,0) to[D,l={普通二极管}] (2.5,0)
        (4,0) to[leD,l={发光二极管}] (6.5,0);
  \node[below] at (0,-0.35) {阳极};
  \node[below] at (2.5,-0.35) {阴极};
\end{circuitikz}
\end{center}

二极管可用于反接保护、钳位、整流和感性负载续流。保护是否有效取决于浪涌
能量、响应速度、位置和回流路径，不能只看原理图上“画了一颗二极管”。

\topic{晶体管：用一个端口控制另一个电流通路}
晶体管是现代数字电路的关键开关。MOSFET 具有栅极、漏极和源极；栅极电压
控制漏源通路。入门时可以把 N 沟道 MOSFET 理解为“栅极足够高时把负载拉向
地”的受控开关，但设计还需检查阈值不是完全导通电压、导通电阻、栅电荷、
体二极管和最大额定值。

\begin{center}
\begin{circuitikz}
  \draw (0,4) node[vcc]{\(V_{\mathrm{DD}}\)}
    to[R,l={负载}] (0,2.4)
    -- (0,1.8);
  \draw (0,1.8) node[nmos,anchor=D] (mos) {};
  \draw (mos.S) -- (0,0) node[ground]{};
  \draw (mos.G) -- (-2,1.0) node[left]{控制输入};
  \node[right] at (0.4,1.0) {N-MOS};
\end{circuitikz}
\end{center}

互补的 P-MOS 与 N-MOS 可以组成 CMOS 反相器：输入低时上管导通、输出被拉高；
输入高时下管导通、输出被拉低。稳定状态下理想静态电流很小，切换时则要给
节点电容充放电。第三章会从这个开关模型进入逻辑门。

\topic{电感、稳压器、晶振和连接器}
\begin{osgridlongtable}{L{0.15\linewidth}L{0.27\linewidth}L{0.48\linewidth}}
部件 & 基本作用 & 读图时必须检查 \\ \hline
\endfirsthead
部件 & 基本作用 & 读图时必须检查 \\ \hline
\endhead
电感 & 储存磁场能量，反对电流突变 & 电感值、饱和电流、直流电阻、磁芯和布局 \\ \hline
LDO & 用线性方式把较高电压稳成较低电压 & 压差、静态电流、功耗、稳定电容与热 \\ \hline
Buck & 用开关、电感和电容高效降压 & 输入范围、开关节点、反馈、补偿、纹波和 EMI \\ \hline
晶体/振荡器 & 提供频率基准 & 频率、负载电容、启动时间、幅度、抖动和使能 \\ \hline
连接器 & 让电源与信号跨机械边界连接 & 方向、pin 1、额定电流、插拔、屏蔽和防呆 \\ \hline
保险丝/TVS & 限制过流或瞬态过压后果 & 触发条件、能量、钳位电压、位置与失效模式 \\ \hline
\end{osgridlongtable}

\section{器件选型、电源转换与时钟}

看到原理图中的元件，应同时追踪：

\[
\text{功能模型}
\rightarrow\text{电气参数}
\rightarrow\text{具体料号}
\rightarrow\text{封装与方向}
\rightarrow\text{PCB/装配约束}.
\]

“这里需要一颗电容”只是功能模型；容量、耐压、介质、温度特性、封装和偏压
降额决定具体器件是否能完成任务。

\topic{怎样读数据手册的边界}
\begin{osgridlongtable}{L{0.21\linewidth}L{0.33\linewidth}L{0.36\linewidth}}
手册区域 & 它回答什么 & 常见误读 \\ \hline
\endfirsthead
手册区域 & 它回答什么 & 常见误读 \\ \hline
\endhead
Absolute Maximum Ratings & 不应超过的损伤边界 & 当成可以长期工作的推荐条件 \\ \hline
Recommended Operating Conditions & 保证功能的工作范围 & 只看典型值，不看最小/最大值 \\ \hline
DC Characteristics & 静态电压、电流、泄漏和阈值 & 忽略测试条件和温度 \\ \hline
AC/Timing Characteristics & 延迟、建立/保持、频率和边沿 & 把典型延迟当最坏保证 \\ \hline
Package/Pinout & 引脚编号、视角、热焊盘和机械尺寸 & 把底视图当顶视图 \\ \hline
Application Circuit & 厂商建议的外围连接与布局 & 不核对型号/版本就机械照抄 \\ \hline
\end{osgridlongtable}

\begin{fromzerobox}
“绝对最大 \(\SI{4}{V}\)”不表示器件适合在 \(\SI{4}{V}\) 工作；它更接近
“越过这里可能损坏，而且接近这里也未承诺正常功能”。设计应落在推荐工作范围，
并为电源误差、纹波、温度和瞬态留裕量。
\end{fromzerobox}

\topic{电阻的实物信息：阻值之外还有功率和温度}
分立电阻可用色环或表面印字表示阻值。贴片电阻“103”常表示
\(10\times10^3=\SI{10}{k\ohm}\)，但 0 欧跳线、精密编码和小封装标记可能
采用不同规则，不能脱离料号猜测。

\begin{center}
\begin{tikzpicture}[font=\footnotesize\sffamily]
  \draw[very thick,BookMuted] (-4,0) -- (-2.5,0);
  \draw[rounded corners=5pt,fill=BookAmber!18,draw=BookAmber,thick]
    (-2.5,-0.65) rectangle (2.5,0.65);
  \draw[very thick,BookMuted] (2.5,0) -- (4,0);
  \foreach \x/\colorname/\label in {
    -1.5/brown/1,
    -0.75/black/0,
    0.0/orange/\(\times10^3\),
    1.55/BookAmber/\(\pm5\%\)} {
    \fill[\colorname] (\x-0.15,-0.65) rectangle (\x+0.15,0.65);
    \node[below=8mm,align=center] at (\x,-0.65) {\label};
  }
  \node[above=6mm] at (0,0.65)
    {教学色环：棕、黑、橙、金 \(\Rightarrow\SI{10}{k\ohm}\pm5\%\)};
\end{tikzpicture}
\end{center}

额定功率通常在规定环境温度下给出，温度升高后需要降额。脉冲承受能力也不等于
连续功率。小型贴片电阻即使平均功率不大，过压、浪涌或局部温升仍可能失效。

\topic{电容类型决定“同样的容量”能否互换}
\begin{osgridlongtable}{L{0.18\linewidth}L{0.25\linewidth}L{0.27\linewidth}L{0.20\linewidth}}
类型 & 特点 & 常见用途 & 主要注意事项 \\ \hline
\endfirsthead
类型 & 特点 & 常见用途 & 主要注意事项 \\ \hline
\endhead
陶瓷 C0G/NP0 & 稳定、损耗低、小容量 & 时钟、滤波、精密模拟 & 容量范围较小 \\ \hline
陶瓷 X7R/X5R & 体积小、容量较大 & 数字电源去耦 & 直流偏压下容量下降 \\ \hline
铝电解 & 大容量、成本低、有极性 & 低频储能和电源输入 & 极性、寿命、ESR、纹波电流 \\ \hline
钽/聚合物 & 容量密度高、ESR 可较低 & 电源轨缓冲 & 浪涌、短路失效和降额 \\ \hline
薄膜 & 稳定、脉冲能力好 & 滤波、定时和高压 & 体积通常较大 \\ \hline
\end{osgridlongtable}

\begin{center}
\begin{tikzpicture}[font=\footnotesize\sffamily]
  \draw[BookBlue,very thick] (0,0) -- (0,2.6);
  \draw[BookBlue,very thick] (0.55,0) -- (0.55,2.6);
  \node[below] at (0.275,-0.1) {无极性电容符号};
  \draw[BookTeal,very thick] (4,0) -- (4,2.6);
  \draw[BookTeal,very thick] (4.6,0.2) arc[start angle=190,end angle=170,radius=7.5];
  \node[BookRed,font=\Large] at (3.65,2.15) {\(+\)};
  \node[below] at (4.3,-0.1) {有极性电容符号};
  \draw[fill=BookAmber!20,draw=BookAmber,rounded corners=2pt]
    (8,-0.1) rectangle (9.6,2.5);
  \draw[fill=BookMuted] (8.0,-0.1) rectangle (8.25,2.5);
  \node[align=center] at (8.8,1.2) {贴片\\电解外观};
  \node[below,align=center] at (8.8,-0.2) {外壳条纹常标负极；\\仍须核对具体料号};
\end{tikzpicture}
\end{center}

\begin{workedexample}
某 \(\SI{10}{\micro F}\)、\(\SI{6.3}{V}\) X5R 陶瓷电容在
\(\SI{3.3}{V}\) 直流偏压下只保留标称容量的 55\%，再考虑 \(-20\%\) 容差，
最坏有效容量近似：

\[
C_{\mathrm{eff}}
=\SI{10}{\micro F}\times0.55\times0.8
=\SI{4.4}{\micro F}.
\]

若稳压器稳定性要求输出端至少 \(\SI{6}{\micro F}\) 有效容量，这颗“标称
\(\SI{10}{\micro F}\)”电容仍不合格。
\end{workedexample}

\topic{二极管族：整流、钳位、发光和续流是不同任务}
二极管的 \(I\)-\(V\) 关系不是固定压降开关。PN 二极管常用近似式：

\[
I_D\approx I_S\left(e^{V_D/(nV_T)}-1\right).
\]

入门计算可用压降近似，但选型还要检查反向耐压、正向电流、浪涌、恢复时间、
结电容和温升。

\begin{osgridlongtable}{L{0.18\linewidth}L{0.32\linewidth}L{0.40\linewidth}}
器件 & 主要特点 & 典型任务 \\ \hline
\endfirsthead
器件 & 主要特点 & 典型任务 \\ \hline
\endhead
普通 PN 二极管 & 成本低，压降和恢复速度依型号变化 & 整流、方向隔离、续流 \\ \hline
肖特基二极管 & 正向压降较低、恢复快、漏电较大 & 电源 OR-ing、低压整流 \\ \hline
TVS & 为吸收短瞬态能量优化 & 接口浪涌和 ESD 钳位的一部分 \\ \hline
稳压管 & 在规定反向区形成近似电压 & 基准或低能量钳位 \\ \hline
LED & 复合发光，颜色对应材料和能隙 & 状态显示、照明、光通信 \\ \hline
\end{osgridlongtable}

\topic{感性负载为什么需要续流路径}
电感电流不能瞬间中断。关断继电器、风扇或电机线圈时，若没有路径继续流动，
电感会产生足以损坏晶体管的高电压。

\begin{center}
\begin{circuitikz}
  \draw (0,4) node[vcc]{\(V_{\mathrm{DD}}\)}
    to[L,l={线圈 \(L\)}] (0,2)
    -- (0,1.6) node[nmos,anchor=D] (mos) {};
  \draw (mos.S) -- (0,0) node[ground]{};
  \draw (mos.G) -- (-1.8,0.8) node[left]{控制};
  \draw (0,2) -- (2.0,2)
    to[D,l_={续流二极管}] (2.0,4)
    -- (0,4);
  \draw[->,BookTeal,thick] (2.35,2.4) -- (2.35,3.6)
    node[midway,right]{关断后的电流};
\end{circuitikz}
\end{center}

二极管位置要让关断电流形成短回路；方向应在正常通电时截止、关断时导通。
不同关断速度和能量要求可能使用 TVS 或主动钳位，而不是普通续流二极管。

\topic{BJT 与 MOSFET：两种受控器件不能只叫“开关”}
\begin{osgridlongtable}{L{0.18\linewidth}L{0.32\linewidth}L{0.40\linewidth}}
比较 & BJT & MOSFET \\ \hline
\endfirsthead
比较 & BJT & MOSFET \\ \hline
\endhead
控制直觉 & 基极电流影响集电极电流 & 栅源电压控制沟道 \\ \hline
稳态驱动 & 需要持续基极驱动电流 & 理想栅极稳态电流很小 \\ \hline
动态驱动 & 受存储电荷影响 & 必须给栅电荷充放电 \\ \hline
导通损耗 & 常用 \(V_{CE(sat)}\) 估算 & 常用 \(I^2R_{DS(on)}\) 估算 \\ \hline
常见误区 & 把 \(\beta\) 当固定精确值 & 把 \(V_{GS(th)}\) 当完全导通电压 \\ \hline
\end{osgridlongtable}

逻辑电平 MOSFET 必须在实际栅压下有保证的
\(R_{DS(on)}\)。数据手册的阈值电压只说明微小测试电流开始出现，并不说明
器件已经适合承载额定负载。

\begin{workedexample}
MOSFET 在 \(V_{GS}=\SI{3.3}{V}\) 时保证
\(R_{DS(on)}\le\SI{80}{m\ohm}\)，负载电流 \(\SI{2}{A}\)。导通损耗上界近似

\[
P=I^2R=(\SI{2}{A})^2\times\SI{0.08}{\ohm}
=\SI{0.32}{W}.
\]

还需用封装热阻估算结温，并检查开关期间损耗、体二极管和浪涌。只看
“额定 \(\SI{30}{A}\)”不能证明板上条件安全。
\end{workedexample}

\topic{稳压：LDO 与 Buck 的能量路径不同}
LDO 像受反馈控制的可变压降器件，近似损耗为

\[
P_{\mathrm{loss,LDO}}\approx
(V_{\mathrm{in}}-V_{\mathrm{out}})I_{\mathrm{out}}.
\]

Buck 在开关节点周期性连接输入与地，借助电感和电容把脉冲能量变成较稳定的
低电压，效率通常更高，但增加纹波、EMI、补偿和布局约束。

\begin{center}
\begin{circuitikz}
  \draw (0,2) node[left]{\(V_{\mathrm{in}}\)}
    -- (0.7,2)
    node[draw=BookBlue,fill=BookCodeBg,minimum width=20mm,minimum height=12mm,
         anchor=west] (switch) {受控开关};
  \draw (switch.east) -- (3.5,2)
    to[L,l={\(L\)}] (5.3,2)
    -- (7.0,2) node[right]{\(V_{\mathrm{out}}\)};
  \draw (5.8,2) to[C,l_={\(C_{\mathrm{out}}\)}] (5.8,0) node[ground]{};
  \draw (3.5,2) to[D,l_={续流路径}] (3.5,0) node[ground]{};
  \draw[BookTeal,dashed,->] (7.0,1.5) -- (6.3,1.5)
    -- (6.3,0.5) -- (1.7,0.5) -- (1.7,1.35);
  \node[BookTeal] at (4.0,0.25) {反馈控制占空比，维持输出};
\end{circuitikz}
\end{center}
\diagramnote{这是非同步 Buck 的功能图，省略补偿、驱动、自举、保护和实际
开关器件；不能直接当成可生产电路。}

\begin{workedexample}
用 LDO 把 \(\SI{12}{V}\) 降到 \(\SI{3.3}{V}\)，电流
\(\SI{0.5}{A}\)：

\[
P_{\mathrm{loss}}\approx(12-3.3)\times0.5=\SI{4.35}{W}.
\]

输出功率只有 \(\SI{1.65}{W}\)，理想效率约
\(3.3/12=27.5\%\)。这通常需要 Buck，而不是简单换一颗“更大封装 LDO”。
\end{workedexample}

\topic{时钟源：晶体、振荡器与 PLL 各承担哪一层}
无源晶体需要放大器和匹配负载才能振荡；有源振荡器模块直接输出时钟；PLL
比较参考时钟与反馈相位，生成不同频率并控制长期相位关系。CPU 内部 GHz
时钟通常不是主板晶体直接输出，而是由较低频参考经过 PLL 倍频。

\begin{pdfdiagram}[node distance=7mm and 12mm]
  \node[os hardware,minimum width=24mm] (crystal) {晶体/参考振荡器\\低频稳定基准};
  \node[os block,right=of crystal,minimum width=22mm] (pll) {PLL\\相位比较与倍频};
  \node[os state,right=of pll,minimum width=22mm] (clocktree) {时钟树\\分频、门控、分配};
  \node[os hardware,right=of clocktree,minimum width=24mm] (loads) {CPU/内存/设备\\各自时钟域};
  \draw[os arrow] (crystal) -- node[above,font=\scriptsize]{参考} (pll);
  \draw[os arrow] (pll) -- node[above,font=\scriptsize]{合成} (clocktree);
  \draw[os arrow] (clocktree) -- node[above,font=\scriptsize]{低偏斜分配} (loads);
\end{pdfdiagram}

频率准确度、周期抖动、相位噪声、启动时间和占空比都可能影响系统。示波器上
“看见在跳”并不等于时钟满足接收器规格。


\section{原理图、PCB、测量与安全上电}

\term{原理图}{schematic}回答“哪些引脚应该电气相连”；\term{PCB}{printed
circuit board}回答“这些网络在真实板材上怎样用铜几何实现”；\term{封装}
{footprint/package}连接逻辑符号与实体焊盘。

\begin{pdfdiagram}
  \node[os block] (symbol) {原理图符号\\U1、引脚名/号};
  \node[os state,right=of symbol] (net) {网络表\\谁与谁相连};
  \node[os hardware,right=of net] (foot) {PCB 封装\\焊盘与外形};
  \node[os hardware,right=of foot] (part) {实体元件\\方向与 pin 1};
  \draw[os arrow] (symbol) -- node[above]{生成} (net);
  \draw[os arrow] (net) -- node[above]{布线} (foot);
  \draw[os arrow] (foot) -- node[above]{装配} (part);
\end{pdfdiagram}

\topic{同一芯片的三种视图}
\begin{center}
\begin{tikzpicture}[font=\footnotesize\sffamily,scale=0.87,transform shape]
  \node[draw=BookBlue,minimum width=28mm,minimum height=22mm,align=center]
    (logic) at (0,0) {U1\\SPI Flash};
  \foreach \y/\name in {0.7/CS\_N,0.2/MISO,-0.3/WP\_N,-0.8/GND}
    \draw (-1.4,\y) -- (-2.2,\y) node[left]{\name};
  \foreach \y/\name in {0.7/VCC,0.2/HOLD\_N,-0.3/SCLK,-0.8/MOSI}
    \draw (1.4,\y) -- (2.2,\y) node[right]{\name};
  \node at (0,-1.65) {逻辑符号：按功能排引脚};

  \draw[rounded corners=2pt,draw=BookMuted,fill=BookCodeBg]
    (5.0,-1.15) rectangle (7.8,1.15);
  \draw[fill=BookMuted] (6.2,0.8) circle (0.09);
  \foreach \y/\n in {0.75/1,0.25/2,-0.25/3,-0.75/4}
    \node[os pin] at (4.65,\y) {\n};
  \foreach \y/\n in {0.75/8,0.25/7,-0.25/6,-0.75/5}
    \node[os pin] at (8.15,\y) {\n};
  \node at (6.4,-1.65) {顶视封装：按物理位置编号};

  \draw[draw=BookAmber,fill=white] (10.8,-1.1) rectangle (13.6,1.1);
  \foreach \y in {0.75,0.25,-0.25,-0.75} {
    \draw[fill=BookAmber!30] (10.45,\y-0.12) rectangle (10.8,\y+0.12);
    \draw[fill=BookAmber!30] (13.6,\y-0.12) rectangle (13.95,\y+0.12);
  }
  \node[align=center] at (12.2,0) {实物俯视\\缺口/圆点标 pin 1};
  \node at (12.2,-1.65) {元件图：有方向但不显示内部逻辑};
\end{tikzpicture}
\end{center}

\begin{failurebox}
最危险的常见错误是把底视图当顶视图、把原理图上的上下位置当物理引脚位置，
或让符号 pin number 与 footprint pad number 不一致。软件无论写得多正确，
错误接到 VCC、GND、时钟或数据脚都不会被软件修复。
\end{failurebox}

\topic{怎样从一张陌生原理图开始}
\begin{enumerate}[leftmargin=2.2em]
  \item 找标题栏、版本、页码和层次入口，确认正在看哪一版。
  \item 找电源入口、GND、稳压器、power-good、时钟和复位。
  \item 找核心芯片位号，按数据手册核对每组电源脚和去耦。
  \item 沿网络标签追踪接口，不把跨页同名网络误认为断开。
  \item 对低有效名称（\code{\_N}、\(\overline{\mathrm{RESET}}\)）标出有效电平。
  \item 对每个连接器核对视角、pin 1、插入方向和保留脚。
  \item 最后才检查数据总线、差分对、终端和软件可见功能。
\end{enumerate}

\topic{PCB 不是把原理图导线画成铜线}
原理图表达理想网络关系；PCB 决定真实电流路径、寄生参数、热、机械和制造。
典型多层板可能包括元件层、连续参考平面、电源层和底层布线。每个信号都有
去程和回程；回流通常选择最小阻抗路径，而不总是人凭直流电阻想象的最短几何
距离。

\begin{center}
\begin{tikzpicture}[font=\footnotesize\sffamily]
  \fill[BookBlue!20] (0,3.3) rectangle (13,3.8);
  \fill[BookTeal!18] (0,2.4) rectangle (13,2.9);
  \fill[BookAmber!18] (0,1.5) rectangle (13,2.0);
  \fill[BookBlue!12] (0,0.6) rectangle (13,1.1);
  \node[left] at (0,3.55) {顶层信号};
  \node[left] at (0,2.65) {连续 GND 平面};
  \node[left] at (0,1.75) {电源平面};
  \node[left] at (0,0.85) {底层信号};
  \draw[BookBlue,very thick,->] (1,3.55) .. controls (4,4.25) and (8,3.0) .. (12,3.55);
  \draw[BookTeal,very thick,->] (12,2.65) .. controls (8,2.25) and (4,3.05) .. (1,2.65);
  \node[BookBlue,above] at (6.5,3.7) {信号去程};
  \node[BookTeal,below] at (6.5,2.45) {相邻参考平面上的回程};
  \draw[BookRed,dashed,thick] (6.5,2.35) -- (6.5,2.95);
  \node[BookRed,align=center] at (9.8,1.75)
    {若参考平面被切缝，\\回流必须绕行并扩大环路};
\end{tikzpicture}
\end{center}

\topic{边沿速度比时钟频率更能决定互连难度}
一个 \(\SI{1}{MHz}\) 方波若上升时间只有 \(\SI{1}{ns}\)，其边沿包含远高于
\(\SI{1}{MHz}\) 的频率成分。互连是否要当传输线处理，常与传播延迟相对于
上升时间的比例有关，而不是只看重复频率。

在 FR-4 PCB 上信号传播速度量级约为真空光速的一半到三分之二。若走线传播
时间已经占上升时间不可忽略的一部分，阻抗不连续、分支和终端就可能产生反射。

\topic{差分对传的是两个节点之差}
差分接收器关心

\[
V_{\mathrm{diff}}=V_P-V_N,
\]

同时两根线相对于地的共模电压必须落在允许范围。P/N 应按规定间距、参考平面、
阻抗和长度关系布线。把两根线互换可能反相逻辑；只接其中一根会失去差分接收
和共模抑制能力。

\topic{去耦布局要画出闭合电流环，而不是只数电容}
高频瞬态电流路径是

\[
\text{电容正端}
\rightarrow \text{芯片 VDD}
\rightarrow \text{芯片内部开关}
\rightarrow \text{芯片 GND}
\rightarrow \text{电容负端}.
\]

\begin{center}
\begin{tikzpicture}[font=\footnotesize\sffamily,>=Latex]
  \draw[rounded corners=3pt,draw=BookBlue,fill=BookCodeBg]
    (5.1,0.6) rectangle (9.3,4.2);
  \node[font=\bfseries] at (7.2,2.4) {数字芯片};
  \node[os pin] (vdd) at (5.0,3.3) {VDD};
  \node[os pin] (gnd) at (5.0,1.5) {GND};
  \draw[fill=BookAmber!20,draw=BookAmber,rounded corners=2pt]
    (1.5,1.6) rectangle (3.1,3.2);
  \node at (2.3,2.4) {\(C_{\mathrm{dec}}\)};
  \node[above] at (2.3,3.25) {正端};
  \node[below] at (2.3,1.55) {负端};
  \draw[BookRed,very thick,->] (3.1,3.0) -- (4.0,3.0) -- (4.0,3.3) -- (vdd);
  \draw[BookRed,very thick,->] (vdd) -- (6.0,3.3) -- (6.0,1.5) -- (gnd);
  \draw[BookRed,very thick,->] (gnd) -- (4.2,1.5) -- (4.2,1.8) -- (3.1,1.8);
  \node[BookRed,align=center] at (5.4,0.35) {目标：面积小、阻抗低的瞬态电流环};

  \draw[BookMuted,dashed,thick] (0,4.8) -- (11.5,4.8);
  \node[BookMuted,above] at (5.75,4.8) {较远稳压器补充平均能量，但不替代局部高频回路};
  \node[os hardware] (reg) at (11.0,2.4) {稳压器};
  \draw[os arrow,dashed] (reg.north) |- (9.0,4.8);
\end{tikzpicture}
\end{center}
\diagramnote{图中位置代表 PCB 物理邻近关系；与只画网络连接的原理图不同。}

\topic{从网络表到一颗 SPI Flash 的完整连接检查}
以 8-pin SPI NOR 为例，真正可启动不只需要四根 SPI 信号：

\begin{osgridlongtable}{L{0.16\linewidth}L{0.25\linewidth}L{0.31\linewidth}L{0.18\linewidth}}
引脚组 & 连接目标 & 必须检查 & 失败表现 \\ \hline
\endfirsthead
引脚组 & 连接目标 & 必须检查 & 失败表现 \\ \hline
\endhead
VCC/GND & 正确电源轨和回流 & 电压范围、去耦、上电斜率 & 完全无响应或不稳定 \\ \hline
CS\_N & 控制器片选 & 默认上拉、时序、是否误选 & 总线无事务或多个器件冲突 \\ \hline
SCLK & 控制器时钟 & 模式、频率、边沿质量 & ID/数据随机错误 \\ \hline
MOSI/IO0 & 控制器输出 & 电压域、方向、串阻 & 命令不能被器件识别 \\ \hline
MISO/IO1 & 控制器输入 & 三态、上拉、采样边沿 & 读取全 0、全 1 或错位 \\ \hline
WP\_N/IO2 & 上拉或 quad 数据 & 启动模式和写保护策略 & 无法编程或 quad 模式失败 \\ \hline
HOLD\_N/IO3 & 上拉或 quad 数据 & 不能悬空、配置一致 & 传输被意外暂停 \\ \hline
\end{osgridlongtable}

\begin{center}
\begin{circuitikz}
  \node[draw=BookBlue,fill=BookCodeBg,minimum width=28mm,minimum height=38mm,
        align=center] (ctrl) at (0,2.5) {SoC/控制器\\SPI master};
  \node[draw=BookAmber,fill=white,minimum width=28mm,minimum height=38mm,
        align=center] (flash) at (8.0,2.5) {U1\\SPI NOR};
  \draw[->,BookBlue,thick] ([yshift=13mm]ctrl.east) --
    node[above,font=\scriptsize]{CS\_N} ([yshift=13mm]flash.west);
  \draw[->,BookBlue,thick] ([yshift=6mm]ctrl.east) --
    node[above,font=\scriptsize]{SCLK} ([yshift=6mm]flash.west);
  \draw[->,BookBlue,thick] ([yshift=-1mm]ctrl.east) --
    node[above,font=\scriptsize]{MOSI} ([yshift=-1mm]flash.west);
  \draw[<-,BookTeal,thick] ([yshift=-8mm]ctrl.east) --
    node[below,font=\scriptsize]{MISO} ([yshift=-8mm]flash.west);
  \draw (6.4,5.6) node[vcc]{\(\SI{3.3}{V}\)}
    to[R,l={\(\SI{10}{k\ohm}\)}] (6.4,4.6)
    -- (flash.north west);
  \draw (9.6,5.6) node[vcc]{\(\SI{3.3}{V}\)}
    to[C,l={\(\SI{100}{nF}\)}] (9.6,4.6)
    -- (flash.north east);
  \draw (9.6,4.6) -- (10.8,4.6) node[right]{VCC};
  \draw (9.6,0.4) node[ground]{} -- (9.6,0.6) -- (flash.south east);
  \node[align=center,BookMuted] at (4.0,0.2)
    {还需核对 WP\_N、HOLD\_N、pin 1、\\封装焊盘号和控制器启动 strap};
\end{circuitikz}
\end{center}

\topic{安全上电：限流、温升和 ESD 先于功能测试}
首次上电应把可恢复故障限制在小能量范围。一个可重复流程：

\begin{enumerate}[leftmargin=2.2em]
  \item 断电检查极性、短路、焊桥、器件方向和电源轨对地阻值；
  \item 将台式电源电压设为目标值，电流限制从保守值开始；
  \item 不装关键器件或保持复位，先验证每条电源轨；
  \item 观察恒压/限流状态、输入电流和异常发热；
  \item 按电源树从输入到末端测量，确认 power-good 与复位顺序；
  \item 最后再接调试器、时钟和高速接口，逐层开放功能。
\end{enumerate}

普通实验室低压电路仍可能因短路产生高温、熔铜和电池起火。市电、高能电池、
大电容和无隔离电源不属于“照着入门图试一试”的范围，应使用合格设备和专业
安全流程。

\topic{ESD 为什么会在“人没感觉”时损坏输入}
人体和衣物可积累静电，放电持续时间短但电压很高。芯片内部保护结构只能吸收
规定能量。防静电工作垫、腕带、接地工具、合适包装和接口 TVS 分别降低产生、
传递和进入敏感节点的风险；它们不能互相完全替代。

\topic{按症状沿物理链定位，而不是一次换完所有元件}
\begin{osgridlongtable}{L{0.18\linewidth}L{0.31\linewidth}L{0.41\linewidth}}
症状 & 先建立的证据 & 下一步 \\ \hline
\endfirsthead
症状 & 先建立的证据 & 下一步 \\ \hline
\endhead
电源一接就限流 & 断电电阻、热像/温升、分段电源树 & 找短路或反装，不继续强行升流 \\ \hline
电压正常但无时钟 & 振荡器使能、供电、参考波形 & 核对负载、布局、启动时间和探头影响 \\ \hline
时钟有但复位不释放 & power-good、复位源、RC/supervisor 时序 & 逐个时钟域核对同步释放 \\ \hline
复位释放但无总线活动 & 启动 strap、CPU 取指地址、片选 & 检查地址译码和非易失介质 \\ \hline
低速正常、高速出错 & 边沿、反射、回流、串扰和采样裕量 & 检查阻抗、终端、参考平面与时序 \\ \hline
加探头后故障消失 & 探头电容/地线改变了节点 & 把测量负载纳入模型，缩短地弹簧 \\ \hline
\end{osgridlongtable}

\begin{failurebox}
原理图网络正确不等于板一定正确。封装焊盘号、元件方向、未装/DNP、焊接、
PCB 参考平面、回流、温升和测量方式都可能让同一网络在实物上表现不同。每次
改动只针对一条被证据支持的假设，并保存改动前后的波形和条件。
\end{failurebox}


\topic{从电源按键到复位释放的物理链}
\begin{pdfdiagram}
  \node[os hardware] (source) {外部电源\\电压与电流能力};
  \node[os hardware,right=of source] (protection) {保护/转换\\保险、TVS、MOSFET};
  \node[os hardware,right=of protection] (rails) {稳压电源轨\\核心、I/O、内存};
  \node[os state,below=of rails] (pg) {Power Good\\电压进入允许窗口};
  \node[os hardware,left=of pg] (clock) {参考时钟\\振荡并稳定};
  \node[os state,left=of clock] (reset) {复位逻辑\\保持后再释放};
  \node[os hardware,below=of reset] (cpu) {CPU 开始取指};
  \draw[os arrow] (source) -- (protection);
  \draw[os arrow] (protection) -- (rails);
  \draw[os arrow] (rails) -- (pg);
  \draw[os arrow] (pg) -- (clock);
  \draw[os arrow] (clock) -- (reset);
  \draw[os arrow] (reset) -- (cpu);
\end{pdfdiagram}

这条链解释了为什么“ROM 中已经有正确代码”仍不足以启动。CPU 必须先处于
允许的电压、温度、时钟和复位条件，取指接口才能产生有意义的事务。第四章会
把 CPU 产生的事务继续连接到 ROM、RAM 与设备。
